% TOP_PROBLEM: {"id": 55, "title": "Frankl's Union-Closed Sets Conjecture", "category_id": 2, "category": {"id": 2, "name": "combinatorics", "display_name": "Combinatorics", "description": "Counting problems, graph theory, discrete structures.", "slug": "combinatorics", "order_index": 2, "created_at": "2026-07-31T15:26:25.671Z"}, "status": "open"}
% FIRST_POSTED: 2026-09-27
% !TeX program = lualatex
% ATTEMPT_STATUS: unresolved
\documentclass[11pt]{article}
\usepackage[margin=25mm]{geometry}
\usepackage{fontspec}
\setmainfont{Latin Modern Roman}
\usepackage{amsmath,amssymb,mathtools}
\usepackage{unicode-math}
\setmathfont{Latin Modern Math}
\usepackage{xurl,hyperref}
\hypersetup{colorlinks=true,urlcolor=blue}
\setcounter{secnumdepth}{0}
\setlength{\parindent}{0pt}
\setlength{\parskip}{5pt}
\setlength{\emergencystretch}{3em}
\begin{document}
\section*{\texorpdfstring{Frankl's Union-Closed Sets Conjecture}{Frankl's Union-Closed Sets Conjecture}}\label{55}
\subsection{Definitions and mathematical statement}
A finite family $\mathcal F$ of finite sets is union-closed if $A\cup B\in\mathcal F$ whenever $A,B\in\mathcal F$. The standard nondegeneracy condition is $\bigcup\mathcal F\ne\varnothing$. Frankl's conjecture then asserts
\[
 \exists x\in\bigcup\mathcal F:\qquad
 \#\{A\in\mathcal F:x\in A\}\ge\frac{|\mathcal F|}{2}.
\]
The empty set may itself be a member. \emph{Scope qualification:} the catalogue says ``other than the empty set,'' which is ambiguous for a family. Excluding only the empty family is insufficient, because $\mathcal F=\{\varnothing\}$ is union-closed and has no element at all. The displayed statement explicitly excludes both degenerate cases, as required for the usual conjecture. The family is a set of sets, not a multiset with arbitrary weights.
\par\smallskip\textit{Formulation and concept references:} \href{https://arxiv.org/pdf/1309.3297}{[S1]}.

\subsection{Short English statement}
In every nondegenerate finite union-closed family, must some element occur in at least half of the sets?
\subsection{Research attempt}
\paragraph{Scope, process, and first gap.}
This GPT-6 Astra Ultra attempt, prepared on 27 September 2026, uses injections induced by union, reduction to a separating family, and explicit frequency counts. The family is unweighted and its union is nonempty. The partial arguments are classical and elementary; the separating-family mechanism is attributed to the work discussed in [S1] and to Falgas--Ravry [E1]. They are reconstructed rather than claimed as new.

The survey [S1] predates the later entropy-based positive constant lower bounds, including Gilmer's [E2]. Thus its historical discussion of a constant-fraction weakening cannot be read as current evidence that even that weakening is open. A positive constant fraction is still different from the exact one-half conclusion sought here. Our first gap is obtaining one element of frequency at least $|\mathcal F|/2$ when neither a small member nor the separating-family size criterion proved below applies.

\paragraph{Singletons give an injection proof.}
Let $m=|\mathcal F|$ and suppose $\{x\}\in\mathcal F$. Map every member not containing $x$ to $A\cup\{x\}$. Union-closure puts the image in $\mathcal F$, it contains $x$, and deleting $x$ recovers $A$. The map is injective. Therefore
\[
 |\{A:x\notin A\}|\le|\{A:x\in A\}|,
\]
which is exactly the conjectured bound for $x$. This includes the case that the empty set is a member.

More generally, if a fixed nonempty $C\in\mathcal F$ has size $k$, the map $A\mapsto A\cup C$ has fibers of size at most $2^k$: preimages have the same trace outside $C$, with at most $2^k$ possible traces inside it. Its image consists of members containing all of $C$. Thus every $x\in C$ occurs in at least $m/2^k$ members. This is valid but becomes weak when $k$ grows. Treating the map as injective for general $C$ would discard precisely these collisions.

\paragraph{A two-element member also suffices.}
Suppose $\{a,b\}\in\mathcal F$, and let $n_{00},n_{10},n_{01},n_{11}$ count members according to membership of $a,b$. The map
\[
 A\longmapsto A\cup\{a,b\}
\]
is injective on the $00$ class into the $11$ class, because their traces outside $\{a,b\}$ distinguish the preimages. Hence $n_{11}\ge n_{00}$. Adding the two frequencies gives
\[
 d(a)+d(b)=n_{10}+n_{01}+2n_{11}
 \ge n_{00}+n_{10}+n_{01}+n_{11}=m.
\]
At least one is therefore at least $m/2$. Notice that this proof does not require either singleton to belong to the family, and does not require the two frequencies to be equal.

A counterexample to the full conjecture, if one exists, consequently has no one-element or two-element member. This is a genuine structural exclusion, but it leaves families whose smallest nonempty member is larger.

\paragraph{Removing indistinguishable elements.}
Two elements of the union are twins if every member contains both or neither. If $x,y$ are twins, delete $y$ from every set. The resulting family remains union-closed, and the deletion map is injective on members: two preimages with the same image have the same membership of $x$, which forces the same membership of $y$.

Thus the number of members and the frequencies of the retained elements are unchanged. Repeating gives a separating family, meaning that each two distinct remaining elements are distinguished by some member. At least one element remains, because the original union was nonempty. A frequent retained element is equally frequent in the original family. This reduction changes the ground-set size but not the frequency denominator; counting deleted twins as independent witnesses would be incorrect.

\paragraph{A complete small-family criterion.}
Now suppose $\mathcal F$ separates its $n$ ground elements, ordered so that
\[
 d(x_1)\le d(x_2)\le\cdots\le d(x_n).
\]
For $i<j$, there exists a member containing $x_j$ but not $x_i$. Otherwise every member containing $x_j$ would contain $x_i$, giving $d(x_j)\le d(x_i)$. Together with the ordering this forces equal occurrence sets, contradicting separation.

For each $i<n$, take the union of chosen witnesses containing $x_j$ but not $x_i$, over all $j>i$, and call it $A_i$. Then
\[
 A_i\in\mathcal F,\qquad x_i\notin A_i,\qquad
 \{x_{i+1},\ldots,x_n\}\subseteq A_i.
\]
These $n-1$ members are distinct: if $i<j$, then $x_j\in A_i$ but $x_j\notin A_j$. The full union $U=\bigcup\mathcal F$ is itself in $\mathcal F$ by finite union-closure and differs from every $A_i$. All these $n$ members contain $x_n$. Therefore
\[
 d(x_n)\ge n.
 \tag{411.1}
\]
For $n=1$, the same conclusion follows from the member $U$, without constructing any $A_i$.

Consequently every separating union-closed family with $m\le2n$ satisfies the conjecture. For an arbitrary family the criterion applies with the number of twin classes after reduction. It is not valid to keep the original, possibly inflated number of ground elements after identifying twins.

The witness construction also proves
\[
 \sum_{A\in\mathcal F}|A|
   =\sum_x d(x)\ge n+(n-1)+\cdots+1
   =\frac{n(n+1)}2.
\]
This is additional incidence information, but it is an absolute bound in terms of $n$. When $m$ is much larger than $n$, it does not force an individual frequency to reach $m/2$.

\paragraph{Why a direct extension stalls.}
The witnesses $A_i$ are indexed by elements, not by all members. The construction produces $n$ sets containing $x_n$, and supplies no injection from every set avoiding $x_n$ into that list. Additional witnesses may coincide with old ones. Thus the numerical implication from (411.1) to the half bound uses $m\le2n$ essentially.

Trying to average over unions has a different issue. If $A,B$ are independent uniform members, then $A\cup B$ lies in the family, but is not generally a uniform member. An element of frequency proportion $p$ belongs to the union with probability $2p-p^2$, which is larger than $p$ for $0<p<1$. That increase describes a changed distribution, not a contradiction. Even for the full power set, uniform membership probabilities are $1/2$, while those of the union are $3/4$. Closure controls the support of the union distribution, not its uniformity.

\paragraph{Why arbitrary weights change the assertion.}
Consider the union-closed family
\[
 \{\varnothing,\{a\},\{a,b\}\}.
\]
Unweighted, $a$ occurs in two of its three members. Give these three sets weights $9,1,1$, respectively. Then the weighted occurrence proportions are $2/11$ and $1/11$, both below one-half. This is not a counterexample to Frankl's conjecture: arbitrary weights are excluded by the statement. It shows that replacing cardinalities by a convenient nonuniform measure requires a separate theorem and cannot be assumed during an entropy or averaging argument.

\paragraph{Verification and outcome.}
The exact checks enumerate all set families on ground sets of size at most four, retain the union-closed nondegenerate ones, and verify the injection, two-element-member, twin-reduction, and separating-witness conclusions. They do not extrapolate small-ground-set truth to all families.

\textbf{UNRESOLVED.} The conjecture is proved here for families containing a member of size one or two, and for separating families with at most twice as many members as ground elements. A possible counterexample must survive these reductions. The remaining gap is an argument yielding the half-frequency bound outside these regimes without changing the uniform measure on the family.

\subsection{Sources}
\begin{itemize}
\item[S1]Bruhn, H., Schaudt, O., 'The journey of the union-closed sets conjecture', arXiv:1309.3297 (2013).
\url{https://arxiv.org/pdf/1309.3297}.
\textit{Formulation source.}
\item[E1]V. Falgas--Ravry, \emph{Minimal Weight in Union-Closed Families}, Electronic Journal of Combinatorics 18 (2011), P95.
\url{https://arxiv.org/abs/1101.2589}.
\item[E2]J. Gilmer, \emph{A constant lower bound for the union-closed sets conjecture}, arXiv:2211.09055.
\url{https://arxiv.org/abs/2211.09055}.
\textit{Later constant-fraction progress, distinct from the one-half target; consulted 27 September 2026.}
\end{itemize}
\end{document}
